% Corte mecánico íntegro: parte_iv.tex:409--499.
% Residencia material del ensamblaje sucesor de 102 capítulos.
% Residencia de control tipado: capítulo 90.
\chapter{Invariants de cohérence mathématique et physique des réalisations hélicoïdales}

\section{Invariants mathématiques du prolongement hélicoïdal}

La théorie est réfutée dans le domaine correspondant si l'un des faits suivants se produit :
\begin{enumerate}
\item un tour nonadique efface la profondeur dans le relèvement ;
\item la règle \(g(K+9)=g(K)\) est confondue avec le retour de l'état ;
\item le lecteur radial ne conserve pas les facteurs et les holonomies préfixes ;
\item \cref{espiral:eq:naturalidad-truncamiento} échoue ;
\item l'action de \(\Gamma_9\) ne satisfait pas l'équivariance \cref{espiral:eq:equivarianza-gamma9} ;
\item la retenue composée ne satisfait pas \cref{espiral:eq:cociclo-carry} ;
\item une courbe classique intervient dans la sélection de \(p,\Lambda,C\) ;
\item la loi \cref{espiral:eq:lp-producto} ou l'inversion \cref{espiral:eq:lp-inversion} échoue dans son domaine ;
\item la vis de Pell ne produit pas \cref{espiral:eq:espiral-log-pell} ;
\item on identifie \(\tau\), \(p\) ou \(\kappa_{\mathrm{F}}\) sans application typée.
\end{enumerate}

\section{Conditions constitutives de la réalisation électromagnétique et matérielle}

Une réalisation électromagnétique ou matérielle est réfutée si :
\begin{enumerate}
\item elle ne satisfait pas \(\mathbf E\perp\mathbf B\perp\widehat{\mathbf k}\), \(|E|=c|B|\) ou \(\omega=ck\) ;
\item elle interprète la dilatation de Pell comme une croissance énergétique sans source ;
\item elle confond l'hélicité photonique avec le spin fermionique ;
\item elle ne préserve pas \(J^2=-I\) et \(\star^2=-I\) dans l'opérateur d'entrelacement déclaré ;
\item la signature spirale n'est pas naturelle sur les multisections de l'atlas ;
\item une étiquette physique sélectionne rétrospectivement une cellule ou une courbe ;
\item l'axe \(\varphi\) de la généalogie des constantes est identifié à l'électron matériel sans application dimensionnelle.
\end{enumerate}

\section{Stratification des théorèmes internes et de leurs réalisations physiques}

\begin{longtable}{p{.31\textwidth}p{.22\textwidth}p{.37\textwidth}}
\caption{Statut probatoire des résultats.}
\label{espiral:tab:estatutos-finales}\\
\toprule
Résultat&Statut&Portée\\
\midrule
\endfirsthead
\toprule
Résultat&Statut&Portée\\
\midrule
\endhead
Courbure mixte APP, régimes TRIT, \(\Gamma_9\), hélice nonadique et vis de Pell&\status{théorème interne antérieur}&Exact dans les domaines définis par les constructions précédentes.\\
Enveloppe affine, lecteur par préfixes, naturalité, équivariance et inversion&\status{théorème de cet ouvrage}&Exact pour les troncatures, les voies typées et le secteur réversible ; \(\sd_4\) certifié.\\
Famille \(L_p\) et cinq publications&\status{théorème de cet ouvrage}&Classe la famille puissance--logarithmique, non toutes les courbes.\\
Carte \(p=\tau_++\tau_-\), \(a=\tau_+-\tau_-\)&\status{formalisation interne exacte}&Caractères élémentaires de la catégorie relevée de 2 demi-couronnes ; les cocycles globaux s'obtiennent par somme de voie.\\
Mode électromagnétique circulaire&\status{réalisation physique}&Solution restreinte de Maxwell dans le fibré phase--position.\\
Signature spirale des particules&\status{réalisation conditionnée}&Proposition sur les multisections ; elle reste soumise au carré explicite de constance, de séparation et de naturalité matérielle.\\
Électron central \((5,5)\)&\status{théorème interne antérieur}&Ancrage matériel unique fixé par l'inversion cellulaire.\\
\bottomrule
\end{longtable}

% REMATE_LOCAL_ESPIRALES_CIERRE
\Needspace{30\baselineskip}
\section{Conclusion}

La théorie obtient une réponse unifiée à l'intuition de départ. Une orbite est un feuillet hélicoïdal lorsque la profondeur est conservée ; l'involution bidirectionnelle produit des feuillets conjugués d'ascension et de descente ; les cercles emboîtés sont des sections de niveau ; les cylindres sont des familles de ces sections. Les spirales classiques ne sont pas choisies comme modèles. Elles apparaissent lorsqu'on publie dans \(\R^2\) des caractères radiaux distincts d'une histoire APP--TRIT--TPK.

Le résultat principal n'est pas une nouvelle liste de courbes. C'est une séparation structurelle entre généalogie et apparence, accompagnée d'un lecteur qui conserve exactement l'information que l'apparence peut oublier. La courbe se situe à la fin. La droite réelle, le plan réel et les géométries de Frenet sont des publications. La phase peut revenir tandis que la mémoire avance ; le cercle peut être intrinsèque ou projeté ; la spirale logarithmique peut être l'ombre exacte d'une vis spectrale ; et le même mécanisme peut admettre des réalisations archimédiennes, électromagnétiques et matérielles sans confondre leurs codomaines.

\begin{resultbox}[title={Classification fonctorielle du lecteur radial et stratification de ses réalisations}]
La composition APP--TRIT--TPK génère une classe d'histoires radiales avec holonomie, orientation, retenue, bord et mémoire. Le lecteur enrichi est naturel sous troncature et équivariant sous l'holonomie nonadique. Sa publication archimédienne produit des courbes puissance--logarithmiques et des suspensions hélicoïdales, mais n'est pas fidèle : une même figure peut conserver des généalogies distinctes. Telle est la structure intime dynamique de l'espace que la classification conventionnelle des courbes n'observe qu'après projection.
\end{resultbox}
